\section{Numerical Accuracy Assessment}
\label{sec:numerical_accuracy}

Numerical accuracy must be assessed independently of convergence
to a steady state. A small residual indicates that the discretized
governing equation is satisfied closely, but it does not necessarily
imply that the numerical solution accurately represents the physical
problem.

The spatial accuracy of the numerical solution can be assessed
by comparing it with an exact solution, when one is available.
Otherwise, the observed order of convergence can be estimated by
comparing numerical solutions obtained on successively refined grids.

\subsection{Accuracy Assessment Using an Exact Solution}

When an exact solution is available, the numerical error is
calculated by comparing the numerical and exact solutions.
For a numerical scheme with spatial order $p$, the error in the
asymptotic range can be expressed as
\begin{equation}
E(h)\approx C h^p,
\label{eq:error_order_exact}
\end{equation}
where $h$ is a characteristic grid spacing and $C$ is a
grid-independent constant.

For two grids with spacings $h_1$ and $h_2$, the ratio of
their errors is
\begin{equation}
\frac{E_1}{E_2}
\approx
\left(\frac{h_1}{h_2}\right)^p.
\label{eq:error_ratio_exact}
\end{equation}

When the grid spacing is halved, $h_2=h_1/2$, the observed
spatial order can be estimated as
\begin{equation}
\boxed{
p\approx
\frac{\ln(E_1/E_2)}{\ln 2}.
}
\label{eq:order_exact}
\end{equation}

The error can be quantified using the root-mean-square (RMS)
norm of the difference between the numerical and exact
temperature fields:
\begin{equation}
E=
\sqrt{
\frac{1}{N_{\mathrm{CV}}}
\sum_{i=1}^{n_x}
\sum_{j=1}^{n_y}
\left[
T_{ij}-T_{\mathrm{exact}}(x_i,y_j)
\right]^2
},
\label{eq:exact_rms_error}
\end{equation}
where $N_{\mathrm{CV}}=n_xn_y$ is the total number of
control volumes, and $(x_i,y_j)$ denotes the center of
each control volume.

\subsection{Grid Convergence Without an Exact Solution}

For the present problem, an exact solution for the complete
temperature field subject to the prescribed inlet condition
is not available. Therefore, the numerical error cannot be
calculated directly using Equation~\eqref{eq:exact_rms_error}.
Instead, the observed spatial order of convergence is estimated
by comparing solutions obtained on three successively refined
grids.

Assuming that the numerical solutions approach the same limiting
solution as the grid is refined, the solution on a grid with
characteristic spacing $h$ can be expressed as
\begin{equation}
T_h=T_{\mathrm{exact}}+C h^p,
\label{eq:grid_solution_error}
\end{equation}
where $T_{\mathrm{exact}}$ represents the limiting solution,
$C$ is the leading error coefficient, and $p$ is the observed
spatial order of convergence.

Let $T_{M_1}$, $T_{M_2}$, and $T_{M_3}$ denote the numerical
solutions on three grids whose characteristic spacings satisfy
\begin{equation}
h_2=\frac{h_1}{2},
\qquad
h_3=\frac{h_2}{2}.
\label{eq:grid_refinement_ratio}
\end{equation}

The differences between successive solutions are defined as
\begin{equation}
D_{12}=\left\|T_{M_1}-T_{M_2}\right\|_2,
\qquad
D_{23}=\left\|T_{M_2}-T_{M_3}\right\|_2,
\label{eq:grid_difference_norms}
\end{equation}
where $\|\cdot\|_2$ denotes the discrete RMS norm. Since the
solutions are defined on different cell-centered grids, their
temperature fields must first be interpolated onto common
physical locations before the differences are calculated.

Assuming that the solutions are in the asymptotic convergence
range, the differences between successive grids satisfy
\begin{equation}
\frac{D_{23}}{D_{12}}\approx\frac{1}{2^p}.
\label{eq:grid_convergence_ratio}
\end{equation}

Taking logarithms gives the observed spatial order:
\begin{equation}
\boxed{
p_{\mathrm{obs}}
=
\frac{\ln(D_{12}/D_{23})}{\ln 2}.
}
\label{eq:observed_spatial_order}
\end{equation}

This approach estimates the spatial convergence order without
requiring an exact solution. Its reliability depends on the
solutions exhibiting consistent convergence and the leading
spatial discretization error dominating other numerical errors.
Therefore, all grids must be sufficiently converged to steady
state so that the remaining pseudo-time convergence error does
not significantly affect the differences between solutions.

Finally, the observed order of convergence does not by itself
establish the absolute accuracy of the finest-grid solution.
It indicates how rapidly the numerical solution approaches
its limiting value as the grid is refined.